\chapter{Color qutrit algebra and finite gauge connection}
\label{chap:algebra-color-qutrit}
\index{color!qutrit}\index{Weyl operators}
\index{special unitary group}\index{Wilson action}

The residual functional \(\kappa_3\) of the
\cref{chap:ocupacion-estadistica} yields three affine classes; the subsequent
atlas realizes their balance \(12+12+12\) in the
\cref{prop:color-atlas}. A triple of labels is not yet a
representation. The algebraic step requires the qutrit chart, its
orientation, and its phase character to be declared. Under those hypotheses, the representation and
its infinitesimal algebra are obtained exactly, without using masses or
chromodynamic data.

\section{Weyl pair and trace algebra \(0\)}
\label{sec:weyl-color-qutrit}

Let
\[
 V_{\rm col}=\mathbb C^3,
 \qquad
 \omega=e^{2\pi i/3},
\]
with ordered basis \((e_0,e_1,e_2)\). We define
\[
 Xe_j=e_{j+1\bmod3},
 \qquad
 Ze_j=\omega^j e_j.
\]
We have
\[
 X^3=Z^3=I,
 \qquad
 ZX=\omega XZ.
\]

\begin{theorem}[Weyl Basis and Complex Color Algebra]
\label{thm:weyl-sl3-color}
For
\[
 W_{ab}:=X^aZ^b,
 \qquad
 (a,b)\in(\mathbb Z/3\mathbb Z)^2,
\]
the nine operators \(W_{ab}\) form an orthogonal basis of
\(\operatorname{End}_{\mathbb C}(V_{\rm col})\) with respect to the
Hilbert--Schmidt product. The eight non-identity operators have zero trace and
\[
 \operatorname{span}_{\mathbb C}
 \{W_{ab}:(a,b)\ne(0,0)\}
 =\mathfrak{sl}_3(\mathbb C).
\]
\end{theorem}

\begin{proof}
The Weyl relation makes it possible to reduce any product to the form
\(X^aZ^b\). A direct calculation gives
\[
 \operatorname{tr}(W_{ab}^{\dagger}W_{cd})
 =3\,\delta_{ac}\delta_{bd}.
\]
Therefore, the nine operators are linearly independent; since
\(\dim_{\mathbb C}\operatorname{End}_{\mathbb C}(V_{\rm col})=9\), they form a
basis. The trace of \(X^aZ^b\) vanishes except for
\((a,b)=(0,0)\). The remaining eight thus form a basis of the trace-zero
subspace, whose complex dimension is eight.
\end{proof}

\begin{corollary}[Compact real form]
\label{cor:su3-color-qutrit}
The Involution
\[
 \tau(A):=-A^\dagger
\]
preserves \(\mathfrak{sl}_3(\mathbb C)\), and its fixed-point set is
\[
 \mathfrak{sl}_3(\mathbb C)^\tau
 =\{A:A^\dagger=-A,\ \operatorname{tr}A=0\}
 =\mathfrak{su}(3).
\]
This real form has dimension eight.
\end{corollary}

\begin{proof}
We have
\[
 W_{ab}^{\dagger}=\omega^{ab}W_{-a,-b}.
\]
The eight non-zero indices form four adjacent pairs. For
\[
 \mathcal R=\{(1,0),(0,1),(1,1),(1,2)\},
\]
the operators
\[
 A_u=W_u-W_u^\dagger,
 \qquad
 B_u=i(W_u+W_u^\dagger),
 \qquad u\in\mathcal R,
\]
are anti-Hermitian, traceless, and linearly independent over
\(\mathbb R\). They therefore constitute a real basis of
\(\mathfrak{su}(3)\).
\end{proof}

The residual application
\[
 \kappa_{\rm col}(r,c)
 =\mathbb C e_{\,r-c\bmod3}
\]
associates the three classes of the atlas with the three lines proper of \(Z\). The
operator \(X\) the permutes cyclically and \(Z\) records its phase. This
assignment is exact a time fixed
\[
 V_{\rm col}=\mathbb C[\mathbb F_3],
 \quad
 \text{the delta basis},
 \quad
 j\longmapsto j+1,
 \quad
 j\longmapsto\omega^j.
\]
Admissible affine changes of the chart produce unitarily conjugate
representations. The construction is therefore natural up to
unitary conjugacy; by itself, it selects neither the physical names of the
colors nor a QCD realization.

\section{Connection and Action on a Finite Cellular Complex}
\label{sec:gauge-finito-color}

\begin{construction}[Finite gauge theory relative to the complex]
\label{cons:gauge-finito-color}
Let \(\mathcal K=(V,E,F)\) be a finite oriented cellularization compatible with
the graph of routes. To each oriented edge \(e\) one assigns
\[
 U_e\in SU(3),
 \qquad
 U_{\bar e}=U_e^{-1}.
\]
A transformation \(g:V\to SU(3)\) acts by
\[
 U_e\longmapsto U_e^g
 =g_{t(e)}U_eg_{s(e)}^{-1}.
\]
For each face \(f\), the ordered product
\[
 U_f=\prod_{e\subset\partial f}U_e^{\epsilon(f,e)}
\]
defines its holonomy. With a base vertex \(v_f\) chosen, one has
\[
 U_f^g=g_{v_f}U_fg_{v_f}^{-1}.
\]
\end{construction}

\begin{proposition}[Gauge invariance and covariant transport]
\label{prop:wilson-color-finito}
For \(\beta\ge0\), the action
\[
 S_{\rm col}^{(\mathcal K,\beta)}[U]
 =\frac{\beta}{3}\sum_{f\in F}
 \left(3-\operatorname{Re}\operatorname{tr}U_f\right)
\]
is gauge-invariant and no negativa. If
\(\psi_v\in V_{\rm col}\) transforma as
\(\psi_v\mapsto g_v\psi_v\), then
\[
 \nabla_e\psi:=U_e\psi_{s(e)}-\psi_{t(e)}
\]
satisfies
\[
 \nabla_e^g\psi^g=g_{t(e)}\nabla_e\psi.
\]
\end{proposition}

\begin{proof}
The holonomy changes by conjugation and the trace is cyclic; of there the
invariance. As each \(U_f\) is unitary of dimension three,
\(\operatorname{Re}\operatorname{tr}U_f\le3\), lo that proves the not
negativity. The covariance of \(\nabla_e\) results substituting the two lawes
of transformation.
\end{proof}

\section{Positivity of Class Kernels and Harmonic Decomposition}

The preceding finite connection does not by itself imply reflection
positivity or a Yang--Mills limit. Before formulating such transfers, it is
appropriate to isolate the classical control that is closed. Let \(G\) be a compact
group and \(\widehat G\) its unitary dual. For each
\(\rho\in\widehat G\), let \(d_\rho\) be its dimension and
\(c_\rho(a)\ge0\) coefficients with sufficient decay, relative to the
growth of \(d_\rho\), for the following series to converge
absolutely.

\begin{definition}[Class nucleus]
\label{pf84:YM-07}
It is defined
\[
 K_a(g,h):=
 \sum_{\rho\in\widehat G}
 c_\rho(a)\,
 \Tr\!\bigl(\rho(g)\rho(h)^*\bigr),
 \qquad g,h\in G.
\]
The choice
\[
 c_\rho(t)=d_\rho e^{-tC_2(\rho)},
 \qquad t>0,
\]
provides the example of the heat kernel. The nonnegativity of the
coefficients does not replace the hypothesis of absolute convergence.
\end{definition}

\begin{theorem}[Positivity of Peter--Weyl]
\label{pf84:YM-08}
For \(g_1,\dots,g_n\in G\) and \(z_1,\dots,z_n\in\CC\),
\[
 \sum_{i,j=1}^n
 \overline z_i z_jK_a(g_i,g_j)\ge0.
\]
\end{theorem}

\begin{proof}
The unitarity of the representations and absolute convergence allow
the sum to be reordered:
\[
 \sum_{i,j}\overline z_i z_jK_a(g_i,g_j)
 =
 \sum_{\rho\in\widehat G}
 c_\rho(a)
 \left\lVert
 \sum_j\overline z_j\rho(g_j)
 \right\rVert_{\mathrm{HS}}^2\ge0.
\]
A negative spectral coefficient may destroy this Gram positivity;
therefore the sign forms part of the hypotheses.
\end{proof}

\begin{lemma}[Abelian Poisson Control]
\label{pf84:YM-09}
For \(G=U(1)\), \(\theta\in\RR\), and \(|a|<1\),
\[
 \sum_{n\in\ZZ}a^{|n|}e^{in\theta}
 =\frac{1-a^2}{1-2a\cos\theta+a^2}.
\]
If nonnegative Fourier coefficients are also required,
\(a\) is restricted to \(0\le a<1\).
\end{lemma}

\begin{proof}
The equality results from summing separately the two geometric series with
positive and negative indices and adding the mode \(n=0\). For \(a<0\), the
scalar identity remains valid but no longer constitutes an example of
nonnegative Peter--Weyl coefficients.
\end{proof}

\begin{remark}
The \cref{pf84:YM-07,pf84:YM-08,pf84:YM-09} are classical definitions and results
of harmonic analysis. They provide a positive control for class kernels; they
do not prove Osterwalder--Schrader positivity, a mass gap, confinement, or a
continuum limit, nor do they promote the finite gauge theory of this chapter
to quantum chromodynamics.
\end{remark}

\begin{exactbox}[title={Algebraic closure of color}]
The \cref{thm:weyl-sl3-color,cor:su3-color-qutrit} construct exactly the
chain
\[
\text{ternary residue}\longrightarrow
\text{qutrit}\longrightarrow
\mathfrak{sl}_3(\mathbb C)\longrightarrow\mathfrak{su}(3).
\]
The \cref{cons:gauge-finito-color,prop:wilson-color-finito} prolongs that
chain to a connection and a gauge action on the finite cellularization
\(\mathcal K\). This is the object proved in the chapter. The
comparison with physical chromodynamics is made subsequently and does not redefine the
qutrit, the balance \(12+12+12\), or the constructed representation.
\end{exactbox}
